% !TeX program = xelatex
\documentclass[12pt,a4paper]{article}
\usepackage{fontspec}
\setmainfont{Times New Roman}
\usepackage[english,russian]{babel}
\usepackage{amsmath,amssymb,mathrsfs,cancel}
\usepackage[margin=22mm]{geometry}
\usepackage{graphicx,float}
\usepackage[unicode,hidelinks]{hyperref}
\DeclareMathOperator{\grad}{grad}
\DeclareMathOperator{\Div}{div}
\newcommand{\Rey}{\mathrm{Re}}
\setlength{\emergencystretch}{3em}
\sloppy
\interfootnotelinepenalty=10000
\newcommand{\topic}[1]{\par\medskip\noindent\textbf{#1}\par\nobreak\smallskip}
\renewcommand{\thefootnote}{\arabic{footnote})}
\begin{document}
\section*{Вопрос 11}
Модель двухфазного течения в пренебрежении капиллярной разностью давлений (модель Маскета — Мереса). Опрокидывание волны повышения насыщенности в задаче об одномерном течении (в задаче Бакли — Леверетта). Условия на скачках насыщенности.

\medskip
\noindent\textbf{Источник:} Г. И. Баренблатт, В. М. Ентов, В. М. Рыжик. \textit{Теория нестационарной фильтрации жидкости и газа}. Москва: Недра, 1972. Глава VI, § 2, с. 155--164.
\par\bigskip
\begin{quote}\sloppy
Выведем систему уравнений двухфазной фильтрации в однородной пористой среде без фазовых переходов. При записи закона фильтрации будем предполагать, что в любой точке каждая из фаз находится целиком лишь в одном из двух крайних состояний: а) связном и подвижном и б) несвязном и поэтому неподвижном. В соответствии со сказанным в предыдущем параграфе, в областях медленного термодинамически равновесного течения для тех жидкостей, которые можно считать связными, можно ввести в расчет зависящие только от насыщенности относительные проницаемости $f_i(s)$ и капиллярное давление $p_c(s)$.

Кроме того, не касаясь гистерезисных явлений, будем рассматривать только однонаправленные процессы, которые наиболее интересны для приложений.

В результате сделанных предположений закон фильтрации можно записать в виде:
\begin{equation}\bar u_i=-\frac{k}{\mu_i}f_i(s)\grad p_i\qquad(i=1,2).\tag{VI.2.1}\label{q11:eq-VI-2-1}\end{equation}
\par\nobreak\smallskip{\small Г. И. Баренблатт, В. М. Ентов, В. М. Рыжик. \textit{Теория нестационарной фильтрации жидкости и газа}. Москва: Недра, 1972. Глава VI, § 2, с. 156.}\par
\end{quote}

\begin{quote}\sloppy
Для определенности будем индекс 1 относить к более смачивающей фазе, насыщенность которой равна $s$. В задачах совместной фильтрации нефти или газа и воды более смачивающей фазой обычно является вытесняющая вода.
\par\nobreak\smallskip{\small Г. И. Баренблатт, В. М. Ентов, В. М. Рыжик. \textit{Теория нестационарной фильтрации жидкости и газа}. Москва: Недра, 1972. Глава VI, § 2, с. 156.}\par
\end{quote}

\begin{quote}\sloppy
Если вытесняемая и вытесняющая фазы --- слабосжимаемые капельные жидкости, влиянием сжимаемости на распределение насыщенности можно пренебречь. Действительно, характерное время нестационарного перераспределения давления за счет сжимаемости имеет порядок $t_1=L^2/\varkappa$, где $\varkappa$ --- коэффициент пьезопроводности, $L$ --- характерный размер. Характерное время вытеснения --- порядка $t_2=L/u$, где $u$ --- средняя скорость фильтрации. Обычно скорость фильтрации --- около $10^{-3}$ см/сек, $L\approx10^4\div10^5$ см, а $\varkappa\approx10^4$ см$^2$/сек. Поэтому $t_1/t_2\approx10^{-2}$, откуда видно, что нестационарные процессы упругого перераспределения давления заканчиваются в начале хода вытеснения.

Если жидкости и пористую среду можно считать несжимаемыми, имеем вместо (VI.2.4) и (IV.2.5)
\begin{equation}m\frac{\partial s}{\partial t}+\Div\bar u_1=0,\qquad m\frac{\partial s}{\partial t}-\Div\bar u_2=0.\tag{VI.2.6}\label{q11:eq-VI-2-6}\end{equation}

\par\nobreak\smallskip{\small Г. И. Баренблатт, В. М. Ентов, В. М. Рыжик. \textit{Теория нестационарной фильтрации жидкости и газа}. Москва: Недра, 1972. Глава VI, § 2, с. 157.}\par
\end{quote}

\begin{quote}\sloppy
В систему уравнений и граничные условия входят следующие размерные определяющие параметры: $\Delta p$ --- перепад давления между скважинами или галереями; $L$ --- характерный размер области течения; коэффициенты $k/\mu_1$ и $k/\mu_2$, а также $p_c^0=\alpha\sqrt{m/k}\cos\Theta$. Перейдем в уравнениях (VI.2.1), (VI.2.2), (VI.2.6) к безразмерным переменным:
\[\begin{gathered}
\frac{p_i}{\Delta p}=P_i;\quad\bar U_i=\frac{\bar u_i}{u_0};\quad\left(\bar u_0=\frac{k\Delta p}{\mu_1L}\right);\quad X=\frac{x}{L};\\
Y=\frac{y}{L};\quad Z=\frac{z}{L};\quad\tau=\frac{u_0t}{L}=\frac{k\Delta p}{\mu_1L^2}t;\quad\varepsilon=\frac{p_c^0}{\Delta p}.
\end{gathered}\]

\par\nobreak\smallskip{\small Г. И. Баренблатт, В. М. Ентов, В. М. Рыжик. \textit{Теория нестационарной фильтрации жидкости и газа}. Москва: Недра, 1972. Глава VI, § 2, с. 157.}\par
\end{quote}

\begin{quote}\sloppy
Параметр $\varepsilon=p_c^0/\Delta p$ мал в большинстве интересных для приложений задач, так отмечалось в § 1. Поэтому член $\bar U_c=\varepsilon f_1(s)J'(s)\grad s$ может быть существен лишь в узких областях, где велик градиент насыщенности. Самый факт существования таких областей будет установлен позже, пока же отметим, что их протяженность должна быть мала по сравнению с основной внешней областью именно потому, что в них велик $|\grad s|$, а $s$ --- ограниченная величина. Это позволяет применить для анализа системы (VI.2.7), (VI.2.8) метод асимптотического сращивания [36а], который заключается в использовании разных масштабов при рассмотрении движения в основной области и в узких зонах резкого изменения насыщенности. Подобно этому течение в пограничном слое вязкой жидкости вблизи стенок имеет иной пространственный масштаб, нежели внешнее течение.

Разделим область фильтрации на внешнюю зону, где $U_c$ --- малая величина, и узкую внутреннюю зону или зоны, где вследствие больших значений $|\grad s|$ членом $U_c$ пренебречь нельзя.

Малость параметра $\varepsilon$ приводит к естественной попытке исследовать течение во внешней зоне, разлагая искомое решение в степенной ряд по этому параметру (внешнее разложение). Первый член внешнего разложения мы получим, полагая $\varepsilon=0$. Исследование структуры внутренней зоны (внутреннее разложение) будет рассмотрено в следующем параграфе.
\par\nobreak\smallskip{\small Г. И. Баренблатт, В. М. Ентов, В. М. Рыжик. \textit{Теория нестационарной фильтрации жидкости и газа}. Москва: Недра, 1972. Глава VI, § 2, с. 158.}\par
\end{quote}

\begin{quote}\sloppy
В одномерном случае (плоском, радиальном или сферическом)
уравнения (VI.2.8) с учетом (VI.2.7) при $\varepsilon=0$ запишем, возвращаясь к размерным переменным, в виде:
\begin{equation}\begin{gathered}
m\frac{\partial s}{\partial t}-\frac{k}{\mu_1}\frac{1}{x^{\nu-1}}\frac{\partial}{\partial x}\left(x^{\nu-1}f_1(s)\frac{\partial p}{\partial x}\right)=0;\\
m\frac{\partial s}{\partial t}+\frac{k}{\mu_2}\frac{1}{x^{\nu-1}}\frac{\partial}{\partial x}\left(x^{\nu-1}f_2(s)\frac{\partial p}{\partial x}\right)=0\qquad(p_1=p_2=p)
\end{gathered}\tag{VI.2.9}\label{q11:eq-VI-2-9}\end{equation}
($\nu=1,2,3$ соответственно для линейного, радиального или сферического течения).

Вычитая второе уравнение (VI.2.9) из первого и интегрируя, получим
\begin{equation}[f_1(s)+\mu_0f_2(s)]\frac{\partial p}{\partial x}=\frac{C(t)}{x^{\nu-1}}\qquad\left(\mu_0=\frac{\mu_1}{\mu_2}\right).\tag{VI.2.10}\label{q11:eq-VI-2-10}\end{equation}
Это равенство выражает постоянство суммарного расхода вдоль трубки тока в силу несжимаемости жидкостей. Определяя $\partial p/\partial x$ из (VI.2.10) и подставляя в любое из уравнений (VI.2.9), получим одно уравнение для $s$:
\begin{equation}\frac{\partial s}{\partial t}+\frac{q(t)}{m}\frac{F'(s)}{x^{\nu-1}}\frac{\partial s}{\partial x}=0,\tag{VI.2.11}\label{q11:eq-VI-2-11}\end{equation}
где $q(t)=-\dfrac{k}{\mu_1}C(t)$ --- суммарный расход жидкости через трубку тока, а
\[F(s)=\frac{f_1(s)}{f_1(s)+\mu_0f_2(s)}.\]
Функция $F(s)$ равна отношению скорости фильтрации (или расхода) вытесняющей фазы к суммарной скорости фильтрации (или к суммарному расходу). Функцию $F(s)$ принято называть функцией распределения фаз.

Введем новые независимые переменные: $Q=\int_{t_0}^t\dfrac{q(t)}{m}\,dt$, $W=\dfrac{x^\nu}{\nu}$.

Величину $W$ можно рассматривать как объем трубки тока между сечениями $x_0$ и $x$ (в частности, при $\nu=1$ и $q(t)=\mathrm{const}$ переменная $Q$ пропорциональна времени, а $W=x$). Тогда вместо (VI.2.11) имеем
\begin{equation}\frac{\partial s}{\partial Q}+F'(s)\frac{\partial s}{\partial W}=0.\tag{VI.2.12}\label{q11:eq-VI-2-12}\end{equation}
Этому уравнению в частных производных первого порядка соответствует следующая система характеристических уравнений:
\begin{equation}\frac{dQ}{1}=\frac{dW}{F'(s)}=\frac{ds}{0}.\tag{VI.2.13}\label{q11:eq-VI-2-13}\end{equation}
Общее решение системы (VI.2.13) представляется в виде:
\begin{equation}\begin{gathered}s=C_1;\\W=QF'(s)+C_2.\end{gathered}\tag{VI.2.14}\label{q11:eq-VI-2-14}\end{equation}
Таким образом, вдоль характеристик уравнения (VI.2.12) $s=\mathrm{const}$, и в плоскости $(W,Q)$ характеристики представляют собой прямые линии. Физически это означает, что каждое значение насыщенности $s$ распространяется со скоростью $dW/dQ$, пропорциональной $F'(s)$. В случае движения плоскими волнами $\nu=1$, а $\dfrac{dW}{dQ}$ есть истинная скорость распространения данного значения насыщенности.

На основании равенств (VI.2.14) общее решение уравнения (VI.2.12) можно формально записать в виде:
\begin{equation}W=QF'(s)+W_0(s),\tag{VI.2.15}\label{q11:eq-VI-2-15}\end{equation}
где функция $W_0(s)$ соответствует начальному распределению насыщенности (при $Q=0$, т. е. $t=t_0$).
\par\nobreak\smallskip{\small Г. И. Баренблатт, В. М. Ентов, В. М. Рыжик. \textit{Теория нестационарной фильтрации жидкости и газа}. Москва: Недра, 1972. Глава VI, § 2, с. 158--160.}\par
\end{quote}

\begin{quote}\sloppy
Вернемся к анализу общего вида решения (VI.2.15).

На рис. VI.7, \textit{а, б} изображены типичные кривые $F(s)$ и $F'(s)$. Функция $F'(s)$ имеет максимум в некоторой точке $s_m$. Поэтому в соответствии с формулой (VI.2.15) два разных значения насыщенности могут иметь одинаковую скорость распространения. В связи с этим получаемая по формуле (VI.2.15) зависимость насыщенности от $W$
может стать неоднозначной --- большие значения «обгоняют» меньшие, как это показано на рис. VI.8.

Неоднозначность формального решения, получаемого из (VI.2.15), означает, что непрерывных решений задачи о вытеснении при заданном начальном условии не существует. Чтобы получить решение, имеющее физический смысл, необходимо вводить разрывы (скачки) насыщенности, т. е. поверхности, на которых значение насыщенности меняется скачком. Такие скачки могут существовать и в начальных условиях.

С ростом $Q$ начальное гладкое распределение насыщенности деформируется, как это видно на рис. VI.8. В некоторый момент (т. е. при некотором $Q$) касательная к кривой $s(W)$ становится вертикальной. Начиная с этого момента, возникает и распространяется скачок насыщенности. Положение скачков насыщенности заранее неизвестно и должно быть найдено в зависимости от времени из решения задачи. На скачках (поверхностях разрыва) должны выполняться условия непрерывности давления и сохранения массы каждой из движущихся фаз.

Выведем условия на скачках для общего случая неодномерного двухфазного течения. Первое из этих условий записывается в виде:
\begin{equation}p^{(1)}=p^{(2)}.\tag{VI.2.19}\label{q11:eq-VI-2-19}\end{equation}
\par\nobreak\smallskip{\small Г. И. Баренблатт, В. М. Ентов, В. М. Рыжик. \textit{Теория нестационарной фильтрации жидкости и газа}. Москва: Недра, 1972. Глава VI, § 2, с. 160--161.}\par
\end{quote}

\begin{quote}\sloppy
Рассмотрим теперь условия сохранения массы каждой из фаз при прохождении поверхности разрыва (скачка) через некоторый элемент объема пористой среды (рис. VI.9), вырезанный по нормали к поверхности разрыва. В силу непрерывности давления сжимаемость
жидкостей для условий на скачке несущественна. Применим к этому элементу уравнение сохранения массы (VI.2.3). Имеем
\[\frac{d}{dt}\left(\int_Tms\,d\tau\right)=mV_n(s^{(1)}-s^{(2)})A+o(T),\]
где $A$ --- площадь поперечного сечения элемента (индексом 1 сверху обозначены величины за скачком, индексом 2 --- перед скачком); $V_n$ --- скорость перемещения скачка по нормали к нему; $o(T)$ --- величина, стремящаяся к нулю быстрее, чем $T$.


Поток жидкости через сечения, параллельные поверхности разрыва, равен $(u_{1n}^{(1)}-u_{1n}^{(2)})A$ ($u_{1n}$ --- проекция скорости фильтрации первой фазы на нормаль к поверхности скачка). Поток, связанный с касательной составляющей скорости, исчезающе мал по сравнению с $u_{1n}A$ при стремлении $\Delta n$ к нулю в силу условия (VI.2.19).

Условие сохранения массы первой жидкости примет тогда вид:
\begin{equation}mV_n(s^{(1)}-s^{(2)})=u_{1n}^{(1)}-u_{1n}^{(2)}.\tag{VI.2.20}\label{q11:eq-VI-2-20}\end{equation}
Равенство (VI.2.20) можно переписать в виде:
\begin{equation}V_n=\frac{u_{1n}^{(1)}-u_{1n}^{(2)}}{m(s^{(1)}-s^{(2)})}.\tag{VI.2.21}\label{q11:eq-VI-2-21}\end{equation}
Условие сохранения массы второй жидкости также сводится к выражению (VI.2.20), поскольку
\begin{equation}u_{1n}^{(1)}+u_{2n}^{(1)}=u_{1n}^{(2)}+u_{2n}^{(2)}.\tag{VI.2.22}\label{q11:eq-VI-2-22}\end{equation}
\par\nobreak\smallskip{\small Г. И. Баренблатт, В. М. Ентов, В. М. Рыжик. \textit{Теория нестационарной фильтрации жидкости и газа}. Москва: Недра, 1972. Глава VI, § 2, с. 161--162.}\par
\end{quote}

\begin{quote}\sloppy
Тогда получим следующее выражение для скорости скачка $V_n=V=dx/dt$:
\begin{equation}\frac{dx}{dt}=\frac{q(t)}{mx^{\nu-1}}\frac{F(s^{(1)})-F(s^{(2)})}{s^{(1)}-s^{(2)}}.\tag{VI.2.24}\label{q11:eq-VI-2-24}\end{equation}
Переходя к переменным $W$ и $Q$, получим
\begin{equation}\frac{dW}{dQ}=\frac{F(s^{(1)})-F(s^{(2)})}{s^{(1)}-s^{(2)}}.\tag{VI.2.25}\label{q11:eq-VI-2-25}\end{equation}
\par\nobreak\smallskip{\small Г. И. Баренблатт, В. М. Ентов, В. М. Рыжик. \textit{Теория нестационарной фильтрации жидкости и газа}. Москва: Недра, 1972. Глава VI, § 2, с. 163.}\par
\end{quote}

\end{document}
